% paper3_discussion.tex -- Discussion and Conclusion of Paper 3, input by
% papers/paper3_intercropping.tex.
\section{Discussion}

A total land equivalent ratio above one at every site shows that the two
crops used the available resources more completely together than apart. The
partial ratios indicate how the advantage arose. Maize, the taller and more
competitive crop, retained most of its sole-crop yield, while cowpea lost a
larger share of its yield to shading and competition yet still contributed
enough to lift the total well above one. This asymmetry is typical of
cereal--legume mixtures, in which complementarity in the use of light,
water and nitrogen allows the combined yield to exceed what either crop could
produce on the same land alone \parencite{brooker2015}.

The advantage was largest at the humid tropical site and smallest at the
temperate site. Where the growing season is long and water is not limiting,
the two crops can occupy the canopy for longer and use more of the resources
that a single crop would leave unused. The comparison must nonetheless be read
with care: the land equivalent ratio depends on the sole-crop yields used to
standardise it, and on the relative proportions of the two crops in the
intercrop \parencite{mead1980}. Both were held constant here, so the ratios
can be compared between sites but not with trials that used other row
arrangements.

The land equivalent ratio measures the efficiency of land use, not the value
of the harvest. Because maize and cowpea grain differ in price and in their
role in the diet, a farmer may prefer an arrangement with a slightly lower
ratio but a larger share of the more valuable crop. The nitrogen that cowpea
leaves in the soil, a benefit to the following cereal \parencite{vanlauwe2010},
is also not captured by a ratio calculated within a single season.

\section{Conclusion}

Maize--cowpea intercropping gave a yield advantage at all three sites, with
mean land equivalent ratios of \res{p3-ler-S1}, \res{p3-ler-S2} and
1.32. Maize contributed the larger share of the advantage at every
site. The ratio should be interpreted together with the value of each crop and
the effect of the legume on the following season. The yields are synthetic
and illustrate the calculation only.
